\scanpage{621}
\mspage{9}
il suffit de voir qu'on a

\[
u_!(\Phi)\subset u_!(\Phi_0)^{\sim}\tag*{$(*)$}
\]

d'où résultera qu'on a $\Phi'\overset{\mathrm{def}}{=}\big(u_!(\Phi)\big)^{\sim}\subset u_!(\Phi_0)^{\sim}$.
Mais on a

\[
u_!(\Phi_0)^{\sim}=\big(u_!(\Phi_0)_*\big)^*\quad \deleted{\text{et }u_!(\Phi_0}
\]

et

\[
u_!(\Phi_0)_*=(u^{\unsure{*|\alpha}})^{-1}(\underbrace{\Phi_{0*}}_{\Psi})=u^{*\, -1}(\Psi)\overset{\mathrm{def}}{=}\Psi',
\]

donc

\[
u_!(\Phi_0)^{\sim}=\Psi'{}^*=\Phi'=u_!(\Phi)^{\sim}
\]

d'où \unsure{(*)}. Cet argument ne fait pas appel
à l'hypothèse ``spéciale'' sur $M$, et peut
s'expliciter en le

% slanted
\margin{devrait passer avant 1.8.}
\uline{Corollaire} 1.9. Sous les conditions générales de
\deleted{\unsure{(}}1.5\deleted{\unsure{)}} ($u_!:M\deleted{\unsure{^{\natural}}}\to M'$ étant un foncteur admettant
un adjoint à droite)\deleted{\unsure{$u^*$}}, soient $\Phi_0,\Phi$ deux
parties de $\mathrm{Fl}\,M$ telles que $\Phi\subset\Phi_0^{\sim}$ (ou ce
qui revient au $=$, $\Phi^{\sim}\subset\Phi_0^{\sim}$). Alors

\[
u_!(\Phi)\subset u_!(\Phi_0)^{\sim}\qquad\text{i.e.}\quad u_!(\Phi)^{\sim}\subset u_!(\Phi_0)^{\sim}.
\]

En effet, cette relation équivaut aussi à
$u_!(\Phi_0)_*\subset u_!(\Phi)_*$, c'est-à-dire (par 1.6)
$(u^*)^{-1}(\Phi_{0*})\subset(u^*)^{-1}(\Phi_*)$, et est \unsure{impliqué|impliquée}
par $\Phi_{0*}\subset\Phi_*$. Cette dernière \unsure{d'autre part}
équivaut à l'hypothèse $\Phi^{\sim}\subset\Phi_0^{\sim}$, \unsure{qed}.

\unsure{!}
% archive sheet number: 620
